\documentclass[10pt,a4paper]{amsart}
\usepackage[margin=19mm]{geometry}
\usepackage{amsmath,amssymb,mathtools}
\usepackage[scr=rsfs,cal=euler]{mathalfa}
\usepackage{xfrac}
\usepackage[unicode,hidelinks,bookmarks=false]{hyperref}
\newcommand{\ie}{i.e.\ }
\newcommand{\cf}{cf.\ }
\newcommand{\resp}{respectively\ }
\newcommand{\Subsection}{\subsection}
\providecommand{\nobreakdash}{\nobreak\hbox{-}}
\setlength{\emergencystretch}{2em}
\multlinegap0pt
\usepackage{amssymb}
\newtheorem{theorem}{Theorem}
\title{On $n$-ary $S$-hyperideals}
\author{Mahdi Anbarloei}
\date{Source: arXiv:2605.19836v1 (CC BY 4.0)}
\begin{document}
\maketitle

\noindent\emph{Source and licence.} On $n$-ary $S$-hyperideals. Version arXiv:2605.19836v1 (CC BY 4.0), distributed by arXiv under the Creative Commons Attribution 4.0 International licence, http://creativecommons.org/licenses/by/4.0/, as recorded in the arXiv licence metadata for this exact version and shown on the abstract page https://arxiv.org/abs/2605.19836v1. Full licence notice: \texttt{licenses/arxiv-2605.19836-CC-BY-4.0.txt}.\par\smallskip

\noindent\textit{Editorial scope.} Counted unit: the article's theorem 11 (source0.tex lines 305--307). The article's complete original proof of it is delivered with it (source0.tex lines 308--339, one \texttt{proof} environment of 32 lines). No mathematics is written, repaired or re-derived: the statement and the proof are the article's own lines, in the article's own notation, and no retained line is substituted or re-flowed.  Retained uncounted as the source-local context the proof needs: lines 244--246.  Nothing was omitted from any retained span, so the package carries no omission record.  No retained line contains a citation, so the wrapper carries no bibliography and no bibliography entry.  No retained line contains a figure environment, an image-inclusion command or drawing code, so no image is delivered and the package declares no asset dependency.  Editorial formatting only, in addition: an AMS \texttt{amsart} preamble that restates the article's own declarations and macros used by the retained lines, namely the environments \texttt{theorem} on separate counters, together with the standard packages those lines need.  The retained environments print in this document as Theorem 1 in order of appearance, whereas the article prints the same items with the numbering of its own full document.  The complete native source of the article as distributed in the arXiv e-print for this exact version is bundled unchanged as \texttt{sources/200944/source0.tex}.
\begin{theorem} \label{4}
Let $S$ be an MS of $\mathscr{A}$ containing $1_\mathscr{A}$ and $Q$ be a hyperideal of $\mathscr{A}$. Then $Q^S=\{x \in \mathscr{A} \ \vert \ g(t,x,1_\mathscr{A}^{(n-2)}) \in Q\ \text{for some}  \ \ t \in S\}$   is the smallest $n$-ary $S$-hyperideal of $\mathscr{A}$ containing $Q$.
\end{theorem}

\begin{theorem} \label{11}
Let $S$ be an MS of $\mathscr{A}$ such that $S \subseteq V(\mathscr{A})$. Then, every $n$-ary $S$-hyperideal of $\mathscr{A}$ is $n$-ary semiprime if and only if $S^{-1}\mathscr{A}$ is regular. 
\end{theorem}

\begin{proof}
$\Longrightarrow$ Assume that $\frac{p}{t}$ is an arbitrary element in $S^{-1}\mathscr{A}$. Put $P=\langle g(p^{(n)}) \rangle$. Let $P \cap S \neq \varnothing$. Then, there exists some $s \in S$ such that $g(s, g(p^{(n)}),1_\mathscr{A}^{(n-2)}) \in S$. This implies that\\ 

$\hspace{2.8cm}\frac{p}{t}=\frac{g \big(g(t^{(n)}),g(s, g(p^{(n)}),1_\mathscr{A}^{(n-2)}),p, 1_\mathscr{A}^{(n-3)}\big)}{g \big(g(t^{(n)}),g(s, g(p^{(n)}),1_\mathscr{A}^{(n-2)}),t ,1_\mathscr{A}^{(n-3)}\big)}$

$\hspace{3.1cm}=\frac{g \big(g(p^{(n)}), g(t^{(n)}),g(s,p,1_\mathscr{A}^{(n-2)}), 1_\mathscr{A}^{(n-3)}\big)}{g \big(g(t^{(n)}),g(s, g(p^{(n)}),1_\mathscr{A}^{(n-2)}),t ,1_\mathscr{A}^{(n-3)}\big)}$

$\hspace{3.1cm}=G \big( G(\frac{p}{t}^{(n)}),\frac{g(t^{(n)})}{g(s,g(p^{(n)}),1_\mathscr{A}^{(n-2)})},\frac{g(s, p,1_\mathscr{A}^{(n-2)})}{t},\frac{1_\mathscr{A}}{1_\mathscr{A}}^{(n-3)}\big).$\\


Hence, $\frac{p}{t}$ is a regular element in $\mathscr{A}$ and so $S^{-1}\mathscr{A}$ is regular. Now, let $P \cap S=\varnothing$. By Theorem \ref{4}, $P^S=\{x \in \mathscr{A} \ \vert \ g(t,x,1_\mathscr{A}^{(n-2)}) \in P\ \text{for some}  \ \ t \in S\}$ is an $n$-ary $S$-hyperideal of $\mathscr{A}$ containing $P$. Hence, $P^S$ is an $n$-ary semiprime hyperideal of $\mathscr{A}$ which implies $p \in P^S$ as $g(p^{(n)}) \in P^S$. Therefore, we obtain $g(r,p,1_\mathscr{A}^{(n-2)}) \in P$ for some $r \in \mathscr{A}$ which means $g(r,p,1_\mathscr{A}^{(n-2)})=g(s,g(p^{(n)}),1_\mathscr{A}^{(n-2)})$ for some $s \in \mathscr{A}$.\\

$\hspace{2.8cm}\frac{p}{t}=\frac{g \big(g(t^{(n-1)},1_\mathscr{A}),r,p,1_\mathscr{A}^{(n-3)}\big)}{g \big(g(t^{(n-1)},1_\mathscr{A}),r,t,1_\mathscr{A}^{(n-3)} \big)}$

$\hspace{3.1cm}=\frac{g \big(g(t^{(n-1)},1_\mathscr{A}),g(r,p,1_\mathscr{A}^{(n-2)}),1_\mathscr{A}^{(n-2)}\big)}{g \big(t^{(n)},r,1_ \mathscr{A}^{(n-2)}\big)}$


$\hspace{03.1cm}=\frac{g \big(g(t^{(n-1)},1_\mathscr{A}),g(s,g(p^{(n)}),1_\mathscr{A}^{(n-2)}),1_\mathscr{A}^{(n-2)}\big)}{g \big(g(t^{(n)}),r, 1_\mathscr{A}^{(n-2)}\big)}$


$\hspace{3.1cm}=\frac{g \big(g(t^{(n-1)},s),g(p^{(n)}),1_\mathscr{A}^{(n-2)}\big)}{g \big ( g(t^{(n)}),r, 1_\mathscr{A}^{(n-2)} \big)}$

$\hspace{3.1cm}=G \big(G(\frac{p}{t}^{(n)}), \frac{g(t^{(n-1)},s)}{r},\frac{1_\mathscr{A}}{1_\mathscr{A}}^{(n-2)}\big).$\\




Then, we conclude that $\frac{p}{t}$ is a regular element in $S^{-1}\mathscr{A}$ and so  $S^{-1}\mathscr{A}$ is regular. 

$\Longleftarrow$  Let  $P$ be an $n$-ary $S$-hyperideal of $\mathscr{A}$  and $g(p^{(n)}) \in P$ for $p \in \mathscr{A}$. Then, there exist $\frac{p_1}{t_1}, \ldots, \frac{p_{n-1}}{t_{n-1}} \in S^{-1}\mathscr{A}$ such that $\frac{p}{1_\mathscr{A}}=G \big(G(\frac{p}{1_\mathscr{A}}^{(n)}), \frac{p_1}{t_1}, \ldots, \frac{p_{n-1}}{t_{n-1}} \big)$  
  as $S^{-1}\mathscr{A}$ is regular and $\frac{p}{1_\mathscr{A}} \in S^{-1}\mathscr{A}$. Since $G(\frac{p}{1_\mathscr{A}}^{(n)})=\frac{g(p^{(n)})}{g(1_\mathscr{A}^{(n)})} \in S^{-1}P$, we obtain $\frac{p}{1_\mathscr{A}}  \in S^{-1}P$. This implies that $g(r,p,1_\mathscr{A}^{(n-2)}) \in P$ for some $r \in S$. Then, we conclude that $p \in P$ as $P$ is an $n$-ary $S$-hyperideal of $\mathscr{A}$. Consequently, every $n$-ary $S$-hyperideal of $\mathscr{A}$ is $n$-ary semiprime. 
\end{proof}

\end{document}
